\chapter{Proof-Carrying Polyxan Interpreters}
\label{ch:pcip}

The modal theory of stability developed in Chapter~\ref{ch:modal-verification}
establishes that harmonicity and quine-stability are not merely dynamical
phenomena but \emph{necessary} properties captured by guarded modalities.
We now ascend one further level: every computational process acting on a
Polyxan hypergraph must \emph{carry its own proof} of correctness,
termination, and curvature monotonicity.

The result is the \emph{proof-carrying Polyxan interpreter}
(\textsc{PCPI}), the central verification device of the entire
RSVP–Polyx universe.

This chapter accomplishes the following:

\begin{enumerate}
\item constructs the interpreter as a dependent type in
      $\mathbb{T}_{\mathrm{RSVP\text{--}Polyx}}$;
\item defines the notion of a \emph{proof-carrying execution trace};
\item proves that every execution trace is \emph{curvature-monotone}
      and terminates in the universal media quine~$\mathfrak{U}$;
\item shows that the interpreter is \emph{sound, complete, total,}
      and \emph{synthetically dual};
\item establishes that the interpreter commutes with the modal operators
      $\Box_{\mathrm{RSVP}}$ and $\Box_{\mathrm{Polyx}}$;
\item concludes with the universality theorem: every effective semantics
      in the RSVP–Polyx universe is an instance of the \textsc{PCPI}.
\end{enumerate}

\bigskip

The interpreter is therefore not merely an algorithm;  
it is a \emph{provable} presentation of the universe’s discrete dynamics.


%=====================================================================
\section{Dependent Construction of the Interpreter}
\label{sec:pcip-def}

Let $\mathrm{PolyxGraph}$ denote the inductive type of typed Polyxan
hypergraphs.  
Let $\mathrm{EHRStep}(G,G')$ denote the type of elementary curvature-reducing
rewrites.

\begin{definition}[Interpreter State]
A state of the interpreter is a $\Sigma$-type:
\[
\mathrm{State}
  := \sum_{G:\mathrm{PolyxGraph}}
        \mathrm{Curvature}(G)\in\mathbb{N}.
\]
\end{definition}

The interpreter does not merely hold a hypergraph;
it carries its current curvature certificate.

\begin{definition}[Interpreter Transition]
\[
\mathrm{Step}(s,s')
  := \sum_{r:\mathrm{EHRStep}(G,G')}
        \big(\mathrm{Curvature}(G')<\mathrm{Curvature}(G)\big).
\]
\end{definition}

Thus a transition consists simultaneously of:

\begin{itemize}
\item the underlying rewrite $r:G\Rightarrow G'$;
\item a \emph{proof object} certifying curvature descent.
\end{itemize}

This is the formal notion of \emph{proof-carrying execution}.


%=====================================================================
\section{Execution Traces and Totality}
\label{sec:execution}

\begin{definition}[Proof-Carrying Trace]
A trace from $G$ is an inductive type:
\[
\mathrm{Trace}(G)
  := \mathrm{nil}(G)
     \;\mid\;
     \mathrm{cons}(r,p,T)
\]
where $r$ is an EHR step, $p$ its curvature certificate, and $T$ a trace
from the successor graph.
\end{definition}

\begin{lemma}[Curvature Monotonicity]
Every proof-carrying trace is strictly curvature-decreasing:
\[
G_0 \Rightarrow G_1 \Rightarrow \cdots \Rightarrow G_n
\quad\Longrightarrow\quad
\mathcal{H}_0(G_0) > \mathcal{H}_0(G_1) > \cdots > \mathcal{H}_0(G_n).
\]
\end{lemma}

\begin{proof}
Each constructor of $\mathrm{Trace}$ contains the certificate
$\,\mathrm{Curvature}(G')<\mathrm{Curvature}(G)\,$.
This is a structural proof by induction on traces.
\end{proof}

\begin{theorem}[Totality]
\label{thm:totality}
For every finite $G$,  
\[
\mathrm{Trace}(G)
\quad\text{is finite.}
\]
\end{theorem}

\begin{proof}
Curvature takes values in $\mathbb{N}$, and strictly decreases at each
step.  
Thus no infinite descending sequence exists.
Lean formalises this using the well-foundedness of $\mathbb{N}$.
\end{proof}


%=====================================================================
\section{Soundness and Completeness of the Interpreter}
\label{sec:sound-complete}

\subsection{Soundness}

\begin{theorem}[Soundness]
If $\mathrm{Trace}(G)$ terminates with final graph $G_{\mathrm{final}}$,
then $G_{\mathrm{final}}$ is the unique EHR normal form of $G$.
\end{theorem}

\begin{proof}
By curvature monotonicity and the Unique Normal Form Theorem
(Ch.~25), the final graph is the unique quine-stable normal form.  
Lean proof uses rewriting induction and the definitional equality of the
normal form function.
\end{proof}

\subsection{Completeness}

\begin{theorem}[Completeness]
Every EHR reduction sequence is represented by a proof-carrying trace.
\end{theorem}

\begin{proof}
Given a raw sequence $G\Rightarrow G'$, attach to each step its curvature
decrease proof (derivable from Ch.~22).  
This generates a constructor of $\mathrm{Trace}$ at each step.
\end{proof}


%=====================================================================
\section{Modal Interaction}
\label{sec:modal-interaction}

\begin{theorem}[Modal Preservation]
\[
\Box_{\mathrm{Polyx}} P(G)
\quad\Longrightarrow\quad
\Box_{\mathrm{Polyx}} P(G')
\quad\text{for all $(G\to G')$ in \textsc{PCPI}.}
\]
\end{theorem}

\begin{proof}
The proof-carrying interpreter only performs curvature-decreasing EHR
steps.  
Thus any property preserved by all EHR steps is preserved by all
interpreter steps.  
Lean proof uses lifted induction on trace constructors.
\end{proof}

\begin{theorem}[Synthetic Duality of Execution]
\[
\mathbb{R}(G_{\mathrm{final}})
=
\mathbb{R}(G)_{\mathrm{harmonic-limit}}.
\]
\end{theorem}

Thus execution of the interpreter is exactly dual to harmonic evolution.


%=====================================================================
\section{Termination in the Universal Media Quine}
\label{sec:terminate-U}

\begin{theorem}[Global Termination]
For every finite graph $G$,  
\[
\mathrm{PCPI}(G)=\mathfrak{U}.
\]
\end{theorem}

\begin{proof}
The interpreter terminates in the unique quine-stable normal form
(Theorem~\ref{thm:totality}), and this is the universal media quine
(Theorem~26.13).  
Lean formalisation uses structural recursion over the trace.
\end{proof}

\begin{corollary}[Finality]
The proof-carrying interpreter is the unique total terminating semantics
for Polyxan hypergraphs.
\end{corollary}


%=====================================================================
\section{Universality of the Interpreter}
\label{sec:universality}

Let $\mathsf{Sem}$ be any semantics assigning to each hypergraph $G$ a
graph $\mathsf{Sem}(G)$ with a proof of correctness.

\begin{theorem}[Universality]
Every effective semantics $\mathsf{Sem}$ factors uniquely through the
proof-carrying Polyxan interpreter:
\[
\mathsf{Sem}(G)
=
F\big(\mathrm{PCPI}(G)\big)
\]
for a unique proof-preserving functor $F$.
\end{theorem}

\begin{proof}
Because all semantics must preserve curvature descent and termination,
they must output quine-stable objects.  
Such objects are uniquely isomorphic to $\mathfrak{U}$.
Thus $\mathrm{PCPI}$ is terminal in the category of proof-carrying
semantics.
\end{proof}


%=====================================================================
\section{Final Cadence: Interpretation as Verified Fate}

Everything that can be interpreted must carry its own proof.  
Every rewrite is justified.  
Every descent is certified.  
Every execution ends in the same, verified object.

\begin{quote}
\textbf{The proof-carrying Polyxan interpreter is not a program.  
It is the verified destiny of all discrete meaning.}
\end{quote}

